% GENERATED BY src/python/lmi/tables/tab02_descriptives.py -- DO NOT EDIT
%==============================================================================%
% tab02_descriptives.tex
%
% Purpose : Tabelle 2 der Replikation -- unstandardisierte deskriptive
%           Statistiken der Modellvariablen. Entspricht Table 2 in Kanas &
%           Kosyakova, S. 16 (Bib-Key: Kanas:2023).
% Source  : SOEPremote job 244565, src/stata/remote/descriptives_micro.do
%           Rohlisting: results/transcripts/src_stata_remote_descriptives_micro.do.eml
%           Transkript: results/comparisons/tab02_descriptives.md
% Needs   : booktabs, tabularx (beide bereits in bachelorarbeit.sty geladen)
% Usage   : % GENERATED BY src/python/lmi/tables/tab02_descriptives.py -- DO NOT EDIT
%==============================================================================%
% tab02_descriptives.tex
%
% Purpose : Tabelle 2 der Replikation -- unstandardisierte deskriptive
%           Statistiken der Modellvariablen. Entspricht Table 2 in Kanas &
%           Kosyakova, S. 16 (Bib-Key: Kanas:2023).
% Source  : SOEPremote job 244565, src/stata/remote/descriptives_micro.do
%           Rohlisting: results/transcripts/src_stata_remote_descriptives_micro.do.eml
%           Transkript: results/comparisons/tab02_descriptives.md
% Needs   : booktabs, tabularx (beide bereits in bachelorarbeit.sty geladen)
% Usage   : % GENERATED BY src/python/lmi/tables/tab02_descriptives.py -- DO NOT EDIT
%==============================================================================%
% tab02_descriptives.tex
%
% Purpose : Tabelle 2 der Replikation -- unstandardisierte deskriptive
%           Statistiken der Modellvariablen. Entspricht Table 2 in Kanas &
%           Kosyakova, S. 16 (Bib-Key: Kanas:2023).
% Source  : SOEPremote job 244565, src/stata/remote/descriptives_micro.do
%           Rohlisting: results/transcripts/src_stata_remote_descriptives_micro.do.eml
%           Transkript: results/comparisons/tab02_descriptives.md
% Needs   : booktabs, tabularx (beide bereits in bachelorarbeit.sty geladen)
% Usage   : % GENERATED BY src/python/lmi/tables/tab02_descriptives.py -- DO NOT EDIT
%==============================================================================%
% tab02_descriptives.tex
%
% Purpose : Tabelle 2 der Replikation -- unstandardisierte deskriptive
%           Statistiken der Modellvariablen. Entspricht Table 2 in Kanas &
%           Kosyakova, S. 16 (Bib-Key: Kanas:2023).
% Source  : SOEPremote job 244565, src/stata/remote/descriptives_micro.do
%           Rohlisting: results/transcripts/src_stata_remote_descriptives_micro.do.eml
%           Transkript: results/comparisons/tab02_descriptives.md
% Needs   : booktabs, tabularx (beide bereits in bachelorarbeit.sty geladen)
% Usage   : \input{tables/tab02_descriptives}
% Author  : Emre Suemer, 2026-08-13
%==============================================================================%

\makeatletter
\@ifundefined{NC@rewrite@Z}{%
  \newcolumntype{Z}{>{\raggedright\arraybackslash\hangindent=1.9em\hangafter=1}X}%
}{}
\makeatother

\begin{table}[htbp]
  \centering
  \caption{Unstandardisierte deskriptive Statistiken der Modellvariablen}
  \label{tab:deskriptiv}
  \footnotesize
  \setlength{\tabcolsep}{4pt}%
  \begin{tabularx}{\textwidth}{@{}Zrrr@{}}
    \toprule
    Variable & Mittelwert/Anteil (SD) & N & Wertebereich \\
    \midrule
    In bezahlter Arbeit & 0{,}18 & 5.467 & 0/1 \\
    Deutsche Sprachkenntnisse & 5{,}40 (2{,}95) & 5.463 & 0--12 \\
    Integrationskurs abgeschlossen & 0{,}29 & 5.197 & 0/1 \\
    Sprachzertifikat erworben & 0{,}43 & 5.455 & 0/1 \\
    H\"aufigkeit des Kontakts zu Deutschen & 2{,}73 (1{,}89) & 5.449 & 0--5 \\
    \textbf{Kreisangebot an Integrationskursen, std.} & 0{,}00 (1{,}00) & 5.467 & $-$1{,}57--5{,}18 \\
    Aufenthaltsdauer, in Monaten (/12) & 2{,}42 (1{,}13) & 5.467 & 0{,}00--5{,}83 \\
    Bildungsjahre vor der Migration & 8{,}98 (5{,}25) & 5.129 & 0--27 \\
    Erwerbserfahrung vor der Migration & 0{,}64 & 5.224 & 0/1 \\
    Weiblich & 0{,}37 & 5.467 & 0/1 \\
    Kinder unter 16 Jahren im Haushalt & 0{,}63 & 5.462 & 0/1 \\
    Einreise mit der Familie & 0{,}34 & 5.467 & 0/1 \\
    Unterst\"utzung durch Familie/Verwandte in DE & 0{,}14 & 5.440 & 0/1 \\
    Alter bei der Erstbefragung & 31{,}74 (9{,}81) & 5.467 & 18--64 \\
    Traumatische Erfahrungen & 0{,}56 & 3.492 & 0/1 \\
    Status des Asylantrags &  & 5.467 &  \\
    \quad -- Keine Entscheidung & 0{,}39 &  & 0/1 \\
    \quad -- Anerkannt & 0{,}38 &  & 0/1 \\
    \quad -- Abgelehnt & 0{,}23 &  & 0/1 \\
    Kreisanteil Ausl\"ander:innen, std. & 0{,}00 (1{,}00) & 5.467 & $-$1{,}82--2{,}99 \\
    Kreisbev\"olkerungsdichte, std. & 0{,}00 (1{,}00) & 5.467 & $-$0{,}86--2{,}73 \\
    Kreisarbeitslosenquote, std. & 0{,}00 (1{,}00) & 5.467 & $-$1{,}86--3{,}04 \\
    Kreissorgen \"uber Zuwanderung, std. & 0{,}00 (1{,}00) & 5.467 & $-$1{,}91--4{,}25 \\
    Kreisanteil CDU/CSU-W\"ahler:innen, std. & 0{,}00 (1{,}00) & 5.467 & $-$2{,}23--2{,}87 \\
    Erhebungsstichprobe &  & 5.467 &  \\
    \quad -- M3 & 0{,}32 &  & 0/1 \\
    \quad -- M4 & 0{,}34 &  & 0/1 \\
    \quad -- M5 & 0{,}33 &  & 0/1 \\
    Herkunftsland &  & 5.467 &  \\
    \quad -- Syrien & 0{,}35 &  & 0/1 \\
    \quad -- Afghanistan & 0{,}20 &  & 0/1 \\
    \quad -- Irak & 0{,}17 &  & 0/1 \\
    \quad -- Eritrea & 0{,}05 &  & 0/1 \\
    \quad -- Iran & 0{,}02 &  & 0/1 \\
    \quad -- \"Ubriger Nahost/Nordafrika & 0{,}02 &  & 0/1 \\
    \quad -- Ehemalige UdSSR & 0{,}04 &  & 0/1 \\
    \quad -- Westbalkan & 0{,}04 &  & 0/1 \\
    \quad -- \"Ubriges Afrika & 0{,}07 &  & 0/1 \\
    \quad -- Sonstige & 0{,}04 &  & 0/1 \\
    Erhebungsjahr &  & 5.467 &  \\
    \quad -- 2016 & 0{,}29 &  & 0/1 \\
    \quad -- 2017 & 0{,}33 &  & 0/1 \\
    \quad -- 2018 & 0{,}23 &  & 0/1 \\
    \quad -- 2019 & 0{,}15 &  & 0/1 \\
    \bottomrule
  \end{tabularx}

  \vspace{0.75ex}
  \parbox{\textwidth}{\footnotesize
    \textit{Anmerkungen:} Eigene Berechnung, IAB-BAMF-SOEP-Befragung von
    Gefl\"uchteten 2016--2019, SOEP-Core v36 (Remote Edition), Analysestichprobe
    (5.467 Personenjahre, 2.730 Personen; vgl.
    Tabelle~\ref{tab:stichprobenselektion}). SD = Standardabweichung;
    std.\ = standardisiert. Die schwankenden Fallzahlen spiegeln
    variablenspezifische Item-Nonresponse wider; die Statistiken werden vor der
    multiplen Imputation berechnet (\texttt{mi extract 0}), also
    vollst\"andige-F\"alle-basiert. S\"amtliche 102 in
    \textcite[16]{Kanas:2023} publizierten Werte werden exakt reproduziert
    (40 Mittelwerte, 40 Fallzahlen, 11 Standardabweichungen, 11 Wertebereiche).
    \\[0.4ex]
    \textbf{Abweichung:} Die drei Kategorien des Asylstatus sind
    in der publizierten Tabelle um eine Position rotiert beschriftet
    (dort \glqq Approved 0,39 / Rejected 0,38 / No decision 0,23\grqq).
    Ma\ss{}geblich sind die Wertelabels des Originalcodes
    (\texttt{00013\_data\_clean.do:1541}), wonach Kategorie~1
    \emph{keine Entscheidung} ist. Die Zahlen sind identisch, die Zuordnung oben
    ist die korrigierte.}
\end{table}

% Author  : Emre Suemer, 2026-08-13
%==============================================================================%

\makeatletter
\@ifundefined{NC@rewrite@Z}{%
  \newcolumntype{Z}{>{\raggedright\arraybackslash\hangindent=1.9em\hangafter=1}X}%
}{}
\makeatother

\begin{table}[htbp]
  \centering
  \caption{Unstandardisierte deskriptive Statistiken der Modellvariablen}
  \label{tab:deskriptiv}
  \footnotesize
  \setlength{\tabcolsep}{4pt}%
  \begin{tabularx}{\textwidth}{@{}Zrrr@{}}
    \toprule
    Variable & Mittelwert/Anteil (SD) & N & Wertebereich \\
    \midrule
    In bezahlter Arbeit & 0{,}18 & 5.467 & 0/1 \\
    Deutsche Sprachkenntnisse & 5{,}40 (2{,}95) & 5.463 & 0--12 \\
    Integrationskurs abgeschlossen & 0{,}29 & 5.197 & 0/1 \\
    Sprachzertifikat erworben & 0{,}43 & 5.455 & 0/1 \\
    H\"aufigkeit des Kontakts zu Deutschen & 2{,}73 (1{,}89) & 5.449 & 0--5 \\
    \textbf{Kreisangebot an Integrationskursen, std.} & 0{,}00 (1{,}00) & 5.467 & $-$1{,}57--5{,}18 \\
    Aufenthaltsdauer, in Monaten (/12) & 2{,}42 (1{,}13) & 5.467 & 0{,}00--5{,}83 \\
    Bildungsjahre vor der Migration & 8{,}98 (5{,}25) & 5.129 & 0--27 \\
    Erwerbserfahrung vor der Migration & 0{,}64 & 5.224 & 0/1 \\
    Weiblich & 0{,}37 & 5.467 & 0/1 \\
    Kinder unter 16 Jahren im Haushalt & 0{,}63 & 5.462 & 0/1 \\
    Einreise mit der Familie & 0{,}34 & 5.467 & 0/1 \\
    Unterst\"utzung durch Familie/Verwandte in DE & 0{,}14 & 5.440 & 0/1 \\
    Alter bei der Erstbefragung & 31{,}74 (9{,}81) & 5.467 & 18--64 \\
    Traumatische Erfahrungen & 0{,}56 & 3.492 & 0/1 \\
    Status des Asylantrags &  & 5.467 &  \\
    \quad -- Keine Entscheidung & 0{,}39 &  & 0/1 \\
    \quad -- Anerkannt & 0{,}38 &  & 0/1 \\
    \quad -- Abgelehnt & 0{,}23 &  & 0/1 \\
    Kreisanteil Ausl\"ander:innen, std. & 0{,}00 (1{,}00) & 5.467 & $-$1{,}82--2{,}99 \\
    Kreisbev\"olkerungsdichte, std. & 0{,}00 (1{,}00) & 5.467 & $-$0{,}86--2{,}73 \\
    Kreisarbeitslosenquote, std. & 0{,}00 (1{,}00) & 5.467 & $-$1{,}86--3{,}04 \\
    Kreissorgen \"uber Zuwanderung, std. & 0{,}00 (1{,}00) & 5.467 & $-$1{,}91--4{,}25 \\
    Kreisanteil CDU/CSU-W\"ahler:innen, std. & 0{,}00 (1{,}00) & 5.467 & $-$2{,}23--2{,}87 \\
    Erhebungsstichprobe &  & 5.467 &  \\
    \quad -- M3 & 0{,}32 &  & 0/1 \\
    \quad -- M4 & 0{,}34 &  & 0/1 \\
    \quad -- M5 & 0{,}33 &  & 0/1 \\
    Herkunftsland &  & 5.467 &  \\
    \quad -- Syrien & 0{,}35 &  & 0/1 \\
    \quad -- Afghanistan & 0{,}20 &  & 0/1 \\
    \quad -- Irak & 0{,}17 &  & 0/1 \\
    \quad -- Eritrea & 0{,}05 &  & 0/1 \\
    \quad -- Iran & 0{,}02 &  & 0/1 \\
    \quad -- \"Ubriger Nahost/Nordafrika & 0{,}02 &  & 0/1 \\
    \quad -- Ehemalige UdSSR & 0{,}04 &  & 0/1 \\
    \quad -- Westbalkan & 0{,}04 &  & 0/1 \\
    \quad -- \"Ubriges Afrika & 0{,}07 &  & 0/1 \\
    \quad -- Sonstige & 0{,}04 &  & 0/1 \\
    Erhebungsjahr &  & 5.467 &  \\
    \quad -- 2016 & 0{,}29 &  & 0/1 \\
    \quad -- 2017 & 0{,}33 &  & 0/1 \\
    \quad -- 2018 & 0{,}23 &  & 0/1 \\
    \quad -- 2019 & 0{,}15 &  & 0/1 \\
    \bottomrule
  \end{tabularx}

  \vspace{0.75ex}
  \parbox{\textwidth}{\footnotesize
    \textit{Anmerkungen:} Eigene Berechnung, IAB-BAMF-SOEP-Befragung von
    Gefl\"uchteten 2016--2019, SOEP-Core v36 (Remote Edition), Analysestichprobe
    (5.467 Personenjahre, 2.730 Personen; vgl.
    Tabelle~\ref{tab:stichprobenselektion}). SD = Standardabweichung;
    std.\ = standardisiert. Die schwankenden Fallzahlen spiegeln
    variablenspezifische Item-Nonresponse wider; die Statistiken werden vor der
    multiplen Imputation berechnet (\texttt{mi extract 0}), also
    vollst\"andige-F\"alle-basiert. S\"amtliche 102 in
    \textcite[16]{Kanas:2023} publizierten Werte werden exakt reproduziert
    (40 Mittelwerte, 40 Fallzahlen, 11 Standardabweichungen, 11 Wertebereiche).
    \\[0.4ex]
    \textbf{Abweichung:} Die drei Kategorien des Asylstatus sind
    in der publizierten Tabelle um eine Position rotiert beschriftet
    (dort \glqq Approved 0,39 / Rejected 0,38 / No decision 0,23\grqq).
    Ma\ss{}geblich sind die Wertelabels des Originalcodes
    (\texttt{00013\_data\_clean.do:1541}), wonach Kategorie~1
    \emph{keine Entscheidung} ist. Die Zahlen sind identisch, die Zuordnung oben
    ist die korrigierte.}
\end{table}

% Author  : Emre Suemer, 2026-08-13
%==============================================================================%

\makeatletter
\@ifundefined{NC@rewrite@Z}{%
  \newcolumntype{Z}{>{\raggedright\arraybackslash\hangindent=1.9em\hangafter=1}X}%
}{}
\makeatother

\begin{table}[htbp]
  \centering
  \caption{Unstandardisierte deskriptive Statistiken der Modellvariablen}
  \label{tab:deskriptiv}
  \footnotesize
  \setlength{\tabcolsep}{4pt}%
  \begin{tabularx}{\textwidth}{@{}Zrrr@{}}
    \toprule
    Variable & Mittelwert/Anteil (SD) & N & Wertebereich \\
    \midrule
    In bezahlter Arbeit & 0{,}18 & 5.467 & 0/1 \\
    Deutsche Sprachkenntnisse & 5{,}40 (2{,}95) & 5.463 & 0--12 \\
    Integrationskurs abgeschlossen & 0{,}29 & 5.197 & 0/1 \\
    Sprachzertifikat erworben & 0{,}43 & 5.455 & 0/1 \\
    H\"aufigkeit des Kontakts zu Deutschen & 2{,}73 (1{,}89) & 5.449 & 0--5 \\
    \textbf{Kreisangebot an Integrationskursen, std.} & 0{,}00 (1{,}00) & 5.467 & $-$1{,}57--5{,}18 \\
    Aufenthaltsdauer, in Monaten (/12) & 2{,}42 (1{,}13) & 5.467 & 0{,}00--5{,}83 \\
    Bildungsjahre vor der Migration & 8{,}98 (5{,}25) & 5.129 & 0--27 \\
    Erwerbserfahrung vor der Migration & 0{,}64 & 5.224 & 0/1 \\
    Weiblich & 0{,}37 & 5.467 & 0/1 \\
    Kinder unter 16 Jahren im Haushalt & 0{,}63 & 5.462 & 0/1 \\
    Einreise mit der Familie & 0{,}34 & 5.467 & 0/1 \\
    Unterst\"utzung durch Familie/Verwandte in DE & 0{,}14 & 5.440 & 0/1 \\
    Alter bei der Erstbefragung & 31{,}74 (9{,}81) & 5.467 & 18--64 \\
    Traumatische Erfahrungen & 0{,}56 & 3.492 & 0/1 \\
    Status des Asylantrags &  & 5.467 &  \\
    \quad -- Keine Entscheidung & 0{,}39 &  & 0/1 \\
    \quad -- Anerkannt & 0{,}38 &  & 0/1 \\
    \quad -- Abgelehnt & 0{,}23 &  & 0/1 \\
    Kreisanteil Ausl\"ander:innen, std. & 0{,}00 (1{,}00) & 5.467 & $-$1{,}82--2{,}99 \\
    Kreisbev\"olkerungsdichte, std. & 0{,}00 (1{,}00) & 5.467 & $-$0{,}86--2{,}73 \\
    Kreisarbeitslosenquote, std. & 0{,}00 (1{,}00) & 5.467 & $-$1{,}86--3{,}04 \\
    Kreissorgen \"uber Zuwanderung, std. & 0{,}00 (1{,}00) & 5.467 & $-$1{,}91--4{,}25 \\
    Kreisanteil CDU/CSU-W\"ahler:innen, std. & 0{,}00 (1{,}00) & 5.467 & $-$2{,}23--2{,}87 \\
    Erhebungsstichprobe &  & 5.467 &  \\
    \quad -- M3 & 0{,}32 &  & 0/1 \\
    \quad -- M4 & 0{,}34 &  & 0/1 \\
    \quad -- M5 & 0{,}33 &  & 0/1 \\
    Herkunftsland &  & 5.467 &  \\
    \quad -- Syrien & 0{,}35 &  & 0/1 \\
    \quad -- Afghanistan & 0{,}20 &  & 0/1 \\
    \quad -- Irak & 0{,}17 &  & 0/1 \\
    \quad -- Eritrea & 0{,}05 &  & 0/1 \\
    \quad -- Iran & 0{,}02 &  & 0/1 \\
    \quad -- \"Ubriger Nahost/Nordafrika & 0{,}02 &  & 0/1 \\
    \quad -- Ehemalige UdSSR & 0{,}04 &  & 0/1 \\
    \quad -- Westbalkan & 0{,}04 &  & 0/1 \\
    \quad -- \"Ubriges Afrika & 0{,}07 &  & 0/1 \\
    \quad -- Sonstige & 0{,}04 &  & 0/1 \\
    Erhebungsjahr &  & 5.467 &  \\
    \quad -- 2016 & 0{,}29 &  & 0/1 \\
    \quad -- 2017 & 0{,}33 &  & 0/1 \\
    \quad -- 2018 & 0{,}23 &  & 0/1 \\
    \quad -- 2019 & 0{,}15 &  & 0/1 \\
    \bottomrule
  \end{tabularx}

  \vspace{0.75ex}
  \parbox{\textwidth}{\footnotesize
    \textit{Anmerkungen:} Eigene Berechnung, IAB-BAMF-SOEP-Befragung von
    Gefl\"uchteten 2016--2019, SOEP-Core v36 (Remote Edition), Analysestichprobe
    (5.467 Personenjahre, 2.730 Personen; vgl.
    Tabelle~\ref{tab:stichprobenselektion}). SD = Standardabweichung;
    std.\ = standardisiert. Die schwankenden Fallzahlen spiegeln
    variablenspezifische Item-Nonresponse wider; die Statistiken werden vor der
    multiplen Imputation berechnet (\texttt{mi extract 0}), also
    vollst\"andige-F\"alle-basiert. S\"amtliche 102 in
    \textcite[16]{Kanas:2023} publizierten Werte werden exakt reproduziert
    (40 Mittelwerte, 40 Fallzahlen, 11 Standardabweichungen, 11 Wertebereiche).
    \\[0.4ex]
    \textbf{Abweichung:} Die drei Kategorien des Asylstatus sind
    in der publizierten Tabelle um eine Position rotiert beschriftet
    (dort \glqq Approved 0,39 / Rejected 0,38 / No decision 0,23\grqq).
    Ma\ss{}geblich sind die Wertelabels des Originalcodes
    (\texttt{00013\_data\_clean.do:1541}), wonach Kategorie~1
    \emph{keine Entscheidung} ist. Die Zahlen sind identisch, die Zuordnung oben
    ist die korrigierte.}
\end{table}

% Author  : Emre Suemer, 2026-08-13
%==============================================================================%

\makeatletter
\@ifundefined{NC@rewrite@Z}{%
  \newcolumntype{Z}{>{\raggedright\arraybackslash\hangindent=1.9em\hangafter=1}X}%
}{}
\makeatother

\begin{table}[htbp]
  \centering
  \caption{Unstandardisierte deskriptive Statistiken der Modellvariablen}
  \label{tab:deskriptiv}
  \footnotesize
  \setlength{\tabcolsep}{4pt}%
  \begin{tabularx}{\textwidth}{@{}Zrrr@{}}
    \toprule
    Variable & Mittelwert/Anteil (SD) & N & Wertebereich \\
    \midrule
    In bezahlter Arbeit & 0{,}18 & 5.467 & 0/1 \\
    Deutsche Sprachkenntnisse & 5{,}40 (2{,}95) & 5.463 & 0--12 \\
    Integrationskurs abgeschlossen & 0{,}29 & 5.197 & 0/1 \\
    Sprachzertifikat erworben & 0{,}43 & 5.455 & 0/1 \\
    H\"aufigkeit des Kontakts zu Deutschen & 2{,}73 (1{,}89) & 5.449 & 0--5 \\
    \textbf{Kreisangebot an Integrationskursen, std.} & 0{,}00 (1{,}00) & 5.467 & $-$1{,}57--5{,}18 \\
    Aufenthaltsdauer, in Monaten (/12) & 2{,}42 (1{,}13) & 5.467 & 0{,}00--5{,}83 \\
    Bildungsjahre vor der Migration & 8{,}98 (5{,}25) & 5.129 & 0--27 \\
    Erwerbserfahrung vor der Migration & 0{,}64 & 5.224 & 0/1 \\
    Weiblich & 0{,}37 & 5.467 & 0/1 \\
    Kinder unter 16 Jahren im Haushalt & 0{,}63 & 5.462 & 0/1 \\
    Einreise mit der Familie & 0{,}34 & 5.467 & 0/1 \\
    Unterst\"utzung durch Familie/Verwandte in DE & 0{,}14 & 5.440 & 0/1 \\
    Alter bei der Erstbefragung & 31{,}74 (9{,}81) & 5.467 & 18--64 \\
    Traumatische Erfahrungen & 0{,}56 & 3.492 & 0/1 \\
    Status des Asylantrags &  & 5.467 &  \\
    \quad -- Keine Entscheidung & 0{,}39 &  & 0/1 \\
    \quad -- Anerkannt & 0{,}38 &  & 0/1 \\
    \quad -- Abgelehnt & 0{,}23 &  & 0/1 \\
    Kreisanteil Ausl\"ander:innen, std. & 0{,}00 (1{,}00) & 5.467 & $-$1{,}82--2{,}99 \\
    Kreisbev\"olkerungsdichte, std. & 0{,}00 (1{,}00) & 5.467 & $-$0{,}86--2{,}73 \\
    Kreisarbeitslosenquote, std. & 0{,}00 (1{,}00) & 5.467 & $-$1{,}86--3{,}04 \\
    Kreissorgen \"uber Zuwanderung, std. & 0{,}00 (1{,}00) & 5.467 & $-$1{,}91--4{,}25 \\
    Kreisanteil CDU/CSU-W\"ahler:innen, std. & 0{,}00 (1{,}00) & 5.467 & $-$2{,}23--2{,}87 \\
    Erhebungsstichprobe &  & 5.467 &  \\
    \quad -- M3 & 0{,}32 &  & 0/1 \\
    \quad -- M4 & 0{,}34 &  & 0/1 \\
    \quad -- M5 & 0{,}33 &  & 0/1 \\
    Herkunftsland &  & 5.467 &  \\
    \quad -- Syrien & 0{,}35 &  & 0/1 \\
    \quad -- Afghanistan & 0{,}20 &  & 0/1 \\
    \quad -- Irak & 0{,}17 &  & 0/1 \\
    \quad -- Eritrea & 0{,}05 &  & 0/1 \\
    \quad -- Iran & 0{,}02 &  & 0/1 \\
    \quad -- \"Ubriger Nahost/Nordafrika & 0{,}02 &  & 0/1 \\
    \quad -- Ehemalige UdSSR & 0{,}04 &  & 0/1 \\
    \quad -- Westbalkan & 0{,}04 &  & 0/1 \\
    \quad -- \"Ubriges Afrika & 0{,}07 &  & 0/1 \\
    \quad -- Sonstige & 0{,}04 &  & 0/1 \\
    Erhebungsjahr &  & 5.467 &  \\
    \quad -- 2016 & 0{,}29 &  & 0/1 \\
    \quad -- 2017 & 0{,}33 &  & 0/1 \\
    \quad -- 2018 & 0{,}23 &  & 0/1 \\
    \quad -- 2019 & 0{,}15 &  & 0/1 \\
    \bottomrule
  \end{tabularx}

  \vspace{0.75ex}
  \parbox{\textwidth}{\footnotesize
    \textit{Anmerkungen:} Eigene Berechnung, IAB-BAMF-SOEP-Befragung von
    Gefl\"uchteten 2016--2019, SOEP-Core v36 (Remote Edition), Analysestichprobe
    (5.467 Personenjahre, 2.730 Personen; vgl.
    Tabelle~\ref{tab:stichprobenselektion}). SD = Standardabweichung;
    std.\ = standardisiert. Die schwankenden Fallzahlen spiegeln
    variablenspezifische Item-Nonresponse wider; die Statistiken werden vor der
    multiplen Imputation berechnet (\texttt{mi extract 0}), also
    vollst\"andige-F\"alle-basiert. S\"amtliche 102 in
    \textcite[16]{Kanas:2023} publizierten Werte werden exakt reproduziert
    (40 Mittelwerte, 40 Fallzahlen, 11 Standardabweichungen, 11 Wertebereiche).
    \\[0.4ex]
    \textbf{Abweichung:} Die drei Kategorien des Asylstatus sind
    in der publizierten Tabelle um eine Position rotiert beschriftet
    (dort \glqq Approved 0,39 / Rejected 0,38 / No decision 0,23\grqq).
    Ma\ss{}geblich sind die Wertelabels des Originalcodes
    (\texttt{00013\_data\_clean.do:1541}), wonach Kategorie~1
    \emph{keine Entscheidung} ist. Die Zahlen sind identisch, die Zuordnung oben
    ist die korrigierte.}
\end{table}
